\chapter{Generating Functionals and the One-Particle-Irreducible Effective Action}
\label{ch:volume-ii-generating-functionals-1pi}
\ConventionsMixed

\section*{Notation for Z-Factors in This Part}
\addcontentsline{toc}{section}{Notation for Z-Factors in This Part}

Several unrelated objects traditionally use the letter \(Z\).  The
renormalization part keeps them distinct as follows.
\begin{center}
\renewcommand{\arraystretch}{1.25}
\begin{tabular}{>{\raggedright\arraybackslash}p{0.23\textwidth}|p{0.67\textwidth}}
symbol & meaning \\
\hline
\(Z[J]\) &
source or partition functional.  Its logarithm generates connected
correlators after a regulator and a branch or formal logarithm have been
specified.  This \(Z\) is not a field-strength factor. \\
\(Z_\phi^{\rm pole}\) &
one-particle K\"all\'en--Lehmann/LSZ residue of a chosen interpolating field:
\(\dd\rho_\phi(s)=Z_\phi^{\rm pole}\delta_{m^2}(\dd s)+\cdots\).  Earlier
scattering chapters often write this residue simply as \(Z_\phi\), or as
\(Z\) when only one scalar field is present.  Once the field normalization is
fixed, it is a physical pole coefficient for that interpolating field and has
no renormalization-scale argument. \\
\(Z_{\rm MOM}(\mu)\) &
finite momentum-subtraction field rescaling at Euclidean subtraction scale
\(\mu\), for example
\([\phi]_\mu=Z_{\rm MOM}(\mu)^{-1/2}\phi_R^{\rm ref}\).  It is a coordinate
choice in the renormalized 1PI chart, not the K\"all\'en--Lehmann pole
residue.  Only after analytic continuation to an isolated physical pole and
the imposition of an on-shell residue condition can a finite field factor be
identified with the pole-residue normalization. \\
\(Z_\phi^{\rm chart}(\lambda(\mu),\mu,\Lambda)\) &
total cutoff-dependent bare-to-renormalized field factor in a chosen
renormalization chart.  In the Callan--Symanzik equation this factor carries
the anomalous dimension through
\(\gamma_\phi=\frac12\,\dd\log Z_\phi^{\rm chart}/\dd\log\mu\) at fixed bare
data. \\
\(Z_\phi^{\rm MS}(\lambda_{\rm MS},\epsilon)\) &
minimal-subtraction field renormalization constant.  It contains pole
counterterms in \(\epsilon\) and no finite residue normalization; it is not the
LSZ residue. \\
\(\mathcal Z^I{}_{A_1\cdots A_r}\) &
operator/source mixing distribution for composite insertions.  This is a
local contact-term renormalization matrix and is unrelated to either the
source functional \(Z[J]\) or a one-particle pole residue.
\end{tabular}
\end{center}
When a later formula suppresses a superscript, the surrounding chart specifies
which row of this table is meant.  In particular, the notation \(Z_\phi\) in
LSZ statements means \(Z_\phi^{\rm pole}\), while \(Z_\phi\) inside a
Callan--Symanzik derivation means the chart field factor
\(Z_\phi^{\rm chart}\).

\section{Sources and Connected Functionals}

\begin{definition}[Regulated scalar source functional]
\label{def:regulated-scalar-source-functional}
Fix a quantum field theory with a specified ultraviolet regulator and a
specified prescription for Lorentzian correlation functions.  Let \(\phi\)
denote a real scalar field in this section; the definitions extend to several
bosonic fields by adding indices, and to fermionic fields by replacing ordinary
sources with Grassmann sources.  Since a quantum field is a distribution, a
c-number source is a real test function, for instance a smooth compactly
supported function or a Schwartz function,
\[
  J\colon \R^{1,D-1}\longrightarrow \R
\]
coupled linearly to \(\phi\) through the pairing
\[
  \langle J,\phi\rangle
  =
  \int \dd^Dx\,J(x)\phi(x).
\]
At finite regulator this pairing is replaced by the corresponding finite
sum, finite-mode pairing, or regulated distributional pairing.  The regulated
source functional is
\[
  Z[J]
  =
  \int D_\Lambda^{\rm ref}\phi\,
  \exp\!\left\{
    \ii S_{\Lambda}[\phi]
    +\ii\langle J,\phi\rangle
  \right\}.
\]
The sign in the Lorentzian source exponent must not be compared directly with
a Euclidean positive-measure convention until the source coordinate has been
continued.  With the lower-half-plane convention of
Chapter~\ref{ch:lorentzian-green-functions}, \(x^0=-\ii\tau\) and
\(\dd x^0=-\ii\dd\tau\), the formal chain for scalar insertions is
\[
  \ii S_L[\phi]+\ii\langle J_L,\phi\rangle_L
  \longmapsto
  -S_E[\phi_E]
  +\langle J_L^{\rm W},\phi_E\rangle_E,
  \qquad
  J_L^{\rm W}(\tau,\vec x):=J_L(-\ii\tau,\vec x).
\]
Thus a Euclidean source functional written with a plus source term uses
\(J_E^{+}=J_L^{\rm W}\), while the positive-measure/free-energy convention
used later in this chapter,
\[
  Z_E[J_E]=\int[D\phi]\,
  \exp\{-S_E[\phi]-\langle J_E,\phi\rangle_E\},
\]
uses \(J_E=-J_L^{\rm W}\).  This is a change of source coordinate, not a new
Green function.  The derivative normalizations and the corresponding
connected-log conventions are collected in
Definition~\ref{def:lorentzian-euclidean-source-bridge}.
With the Fourier convention
\[
  \phi(x)=\int\frac{\dd^Dk}{(2\pi)^D}\,
  \ee^{\ii k\cdot x}\,\widetilde\phi(k),
\]
the source pairing is
\[
  \langle J,\phi\rangle
  =
  \int\frac{\dd^Dk}{(2\pi)^D}\,
  \widetilde J(k)\widetilde\phi(-k).
\]
Perturbatively, the source is therefore a one-valent vertex with factor
\(\ii\widetilde J(k)\):
\[
  \begin{tikzpicture}[baseline=-0.35ex,>=Stealth,scale=0.82,
    src/.style={circle,fill=blue!70!black,inner sep=1.9pt},
    qline/.style={line width=0.75pt,-{Latex[length=1.7mm]}}]
    \node[src] (J) at (0,0) {};
    \draw[qline] (1.25,0)--(J);
    \node[blue!70!black] at (0,-0.42) {\(J\)};
    \node at (0.72,0.32) {\(-k\)};
  \end{tikzpicture}
  \quad =\quad
  \ii\widetilde J(k).
\]
Here \(D_\Lambda^{\rm ref}\phi\) is the reference integration density of
Definition~\ref{def:regulated-scalar-integration-conventions} in a
finite-dimensional bosonic scalar regulator, and \(S_{\Lambda}\) is the
regulated action placed in the exponential.  If a chapter instead uses a
Gaussian reference measure, it writes \(\dd\mu_{C_\Lambda}\) and moves the
quadratic kinetic term into that measure.  Other regulators, such as
fermionic, gauge, BV, Lorentzian oscillatory, or dimensional-regularization
prescriptions, must specify their own regulated integration object.  The
notation suppresses \(\Lambda\) after this paragraph.  All functional
derivatives in this chapter are first derivatives of the regulated source
functional; continuum statements are obtained only after the chosen
renormalization prescription gives a distributional limit.
\end{definition}

\begin{definition}[Status of \(\log Z\)]
\label{def:source-functional-log-status}
The symbol \(\log Z[J]\) has one of the following meanings, fixed by the
regulator and limiting prescription.
\begin{enumerate}
  \item \emph{Formal perturbation theory.}  Let \(g\) denote the interaction
  couplings and work in a formal power-series ring in \(g\), with coefficients
  that are regulated source distributions.  If \(Z[0]\) is invertible, write
  \[
    Z[J]=Z[0]\bigl(1+A[J]\bigr),
    \qquad
    A[0]=0 .
  \]
  Then
  \[
    \log Z[J]
    =
    \log Z[0]
    +
    \sum_{n\geq1}\frac{(-1)^{n+1}}{n}A[J]^n
  \]
  is the algebraic logarithm in that completed ring.  The first term is a
  vacuum functional; source derivatives of positive order depend only on the
  second term.  This is the interpretation used by dimensional
  regularization, minimal subtraction, and graph-by-graph perturbative
  constructions unless an actual positive measure or oscillatory integral has
  also been supplied.

  \item \emph{Positive Euclidean construction.}  If \(Z_E[J]\) is defined by a
  positive finite Euclidean measure on a real source domain and
  \(0<Z_E[J]<\infty\), then \(\log Z_E[J]\) is the ordinary real logarithm.
  Convexity and Legendre--Fenchel statements in the later Euclidean section
  use this meaning and require the positivity and finiteness hypotheses stated
  there.

  \item \emph{Lorentzian oscillatory or phase construction.}  If a finite
  Lorentzian regulator gives a complex-valued source functional \(Z_L[J]\),
  then \(\log Z_L[J]\) is defined only on a source neighborhood on which
  \(Z_L[J]\neq0\) and a branch has been chosen from a base source.  A principal
  branch is a possible convention only when the image of that neighborhood
  avoids its branch cut.  Changing the branch on a connected zero-free
  neighborhood changes \(\log Z_L\) by an additive constant in
  \(2\pi\ii\mathbb Z\), hence changes \(W_L=-\ii\log Z_L\) by a source-
  independent constant and leaves connected correlators of positive order
  unchanged.
\end{enumerate}
No statement in this chapter uses an unclassified logarithm.  Every later
Legendre transform or connected-kernel formula is to be read in the
corresponding one of these three senses.
\end{definition}

\begin{definition}[Connected functional and source-dependent field]
\label{def:connected-functional-source-dependent-field}
For the source functional of
Definition~\ref{def:regulated-scalar-source-functional}, define the
Lorentzian connected functional \(W\) by
\[
  Z[J]=\ee^{\ii W[J]},
  \qquad
  W[J]=-\ii\log Z[J],
\]
with the meaning of \(\log Z\) supplied by
Definition~\ref{def:source-functional-log-status}.  In the Lorentzian
complex-valued case this includes the choice of a zero-free source
neighborhood and a logarithm branch on it; in the formal case it is the
algebraic logarithm and no analytic branch is involved.  The constant \(W[0]\)
is the vacuum functional; subtracting it removes vacuum bubbles and leaves all
source derivatives of positive order unchanged.  The functional \(W[J]\), with
this harmless constant suppressed, generates connected Green functions with
the Lorentzian phase factor
\[
  G^{(n)}_{\rm c}(x_1,\ldots,x_n)
  =
  \left.
  \ii^{1-n}
  \frac{\delta^n W[J]}
       {\delta J(x_1)\cdots \delta J(x_n)}
  \right|_{J=0}.
\]
The full Green functions are generated by \(Z[J]\):
\[
  G^{(n)}(x_1,\ldots,x_n)
  =
  \left.
  \frac1{Z[0]}
  \frac{\delta}{\ii\delta J(x_1)}\cdots
  \frac{\delta}{\ii\delta J(x_n)}\,Z[J]
  \right|_{J=0}.
\]
The correlation-ordering prescription is the one built into the path integral.
With the Feynman boundary prescription specified by the \(+\ii0\) deformation
of quadratic denominators, these are time-ordered Lorentzian Green functions.

The source-dependent one-point function is
\[
  \varphi_J(x)
  =
  \frac{\delta W[J]}{\delta J(x)}
  =
  \frac1{Z[J]}
  \int D_\Lambda^{\rm ref}\phi\,
  \phi(x)\,
  \exp\!\left\{
    \ii S_\Lambda[\phi]+\ii\langle J,\phi\rangle
  \right\}.
\]
At \(J=0\), this is the vacuum expectation value of \(\phi(x)\) in the chosen
state and prescription:
\[
  \varphi_0(x)=\vev{\phi(x)} .
\]
\end{definition}

Figure~\ref{fig:volume-ii-source-legendre} records the source-to-field map used
in the Legendre transform: differentiating \(W_\Lambda[J]\) inserts one field,
and solving for \(J\) as a function of \(\varphi\) defines the 1PI coordinate.

\begin{figure}[!ht]
  \centering
  \resizebox{0.95\textwidth}{!}{%
  \begin{tikzpicture}[>=Stealth,x=1cm,y=1cm,every node/.style={font=\small}]
    \tikzset{
      source/.style={circle,fill=blue!70!black,inner sep=1.8pt},
      line/.style={line width=0.9pt},
      blob/.style={circle,draw=black,fill=gray!18,line width=0.9pt,minimum size=1.0cm},
      arr/.style={-{Latex[length=2.0mm]},line width=0.9pt}
    }
    \begin{scope}
      \node[font=\normalsize] at (0,2.3) {source insertions};
      \node[blob] (B) at (0,0) {};
      \foreach \ang/\name in {25/a,95/b,175/c,245/d,315/e} {
        \coordinate (P\name) at (\ang:1.75);
        \draw[line] (B) -- (P\name);
        \node[source] at (P\name) {};
      }
      \node[blue!70!black] at (0,-2.0) {\(\delta/(\ii\delta J)\)};
      \node[align=center] at (0,-2.65) {connected components\\come from \(W=-\ii\log Z\)};
    \end{scope}

    \draw[arr] (2.45,0.15) -- (4.3,0.15);
    \node[above] at (3.35,0.15) {\(\varphi=\delta W/\delta J\)};

    \begin{scope}[shift={(6.4,0)}]
      \draw[->,line width=0.75pt] (-2.0,-1.35)--(2.2,-1.35) node[right] {\(J\)};
      \draw[->,line width=0.75pt] (-1.6,-1.65)--(-1.6,1.55) node[above] {\(W[J]\)};
      \draw[blue!70!black,line width=1.05pt]
        plot[smooth,tension=0.65] coordinates {
          (-1.75,-1.1) (-1.15,-0.15) (-0.45,0.55) (0.35,0.95) (1.35,1.15) (2.0,1.22)
        };
      \draw[green!55!black,line width=0.95pt] (-1.2,-0.25)--(1.65,1.22);
      \node[green!50!black] at (1.45,0.66) {slope \(\varphi\)};
      \node[font=\normalsize] at (0,2.3) {Legendre data};
    \end{scope}

    \draw[arr] (8.95,0.15) -- (10.7,0.15);
    \node[above] at (9.85,0.15) {\(\Gamma=W-\int\varphi J\)};

    \begin{scope}[shift={(12.6,0)}]
      \draw[->,line width=0.75pt] (-1.8,-1.35)--(2.1,-1.35) node[right] {\(\varphi\)};
      \draw[->,line width=0.75pt] (-1.45,-1.65)--(-1.45,1.55) node[above] {\(\Gamma[\varphi]\)};
      \draw[purple!80!black,line width=1.05pt]
        plot[smooth,tension=0.75] coordinates {
          (-1.45,1.2) (-0.85,0.18) (-0.2,-0.58) (0.55,-0.25) (1.25,0.35) (1.95,1.15)
        };
      \node[purple!80!black] at (0.85,1.3) {\(\delta\Gamma/\delta\varphi=-J\)};
      \node[font=\normalsize] at (0,2.3) {1PI functional};
    \end{scope}
  \end{tikzpicture}
  }
  \caption{The logarithm of the source functional generates connected
  diagrams.  The source-dependent field \(\varphi_J\) is the slope of
  \(W[J]\), and the Legendre transform exchanges the source \(J\) for the
  field \(\varphi\).}
  \label{fig:volume-ii-source-legendre}
\end{figure}

\section{Legendre Transform}

\begin{definition}[Local one-particle-irreducible Legendre transform]
\label{def:local-1pi-legendre-transform}
Assume that, in a neighborhood of the source configurations under discussion,
the map
\[
  J\longmapsto \varphi_J=\frac{\delta W[J]}{\delta J}
\]
is locally invertible.  Denote the inverse by \(J_{\varphi}\), so that
\[
  \varphi(x)=
  \left.
  \frac{\delta W[J]}{\delta J(x)}
  \right|_{J=J_{\varphi}} .
\]
The one-particle-irreducible effective action is the Legendre transform
\[
  \Gamma[\varphi]
  =
  W[J_{\varphi}]
  -
  \langle J_{\varphi},\varphi\rangle .
\]
This is a local differentiable Legendre transform in the chosen source chart.
It is the definition used by perturbative 1PI diagrammatics: \(W\) and
\(\Gamma\) may be formal power series, and \(J_\varphi\) is obtained
coefficient by coefficient by inverting
\(\varphi=\delta W/\delta J\).  This formal operation does not require a
positive measure and does not imply convexity.  In particular, in dimensional
regularization with minimal subtraction the regulated functional integral is
not a positive finite measure on field configurations; it is a prescription
for formal amplitudes.  Whenever later perturbative MS or BPHZ discussions use
the 1PI effective action, this local formal 1PI functional is the operative
object unless a positive Euclidean regulator and the convex
Legendre--Fenchel construction are explicitly invoked.
\end{definition}

\begin{proposition}[1PI source equation and inverse two-point kernel]
\label{prop:1pi-source-equation-inverse-kernel}
For the local Legendre transform of
Definition~\ref{def:local-1pi-legendre-transform},
\[
  \frac{\delta\Gamma[\varphi]}{\delta\varphi(x)}
  =
  -J_{\varphi}(x),
\]
and the connected two-point kernel is the inverse of the 1PI two-point kernel
with the Lorentzian sign
\[
  \int \dd^Dz\,
  W^{(2)}(x,z)\Gamma^{(2)}(z,y)
  =
  -\delta^{(D)}(x-y).
\]
\end{proposition}

\begin{proof}
Varying the definition of \(\Gamma\) gives
\[
\begin{aligned}
  \delta\Gamma[\varphi]
  &=
  \int \dd^Dx\,
  \frac{\delta W[J_{\varphi}]}{\delta J_{\varphi}(x)}
  \delta J_{\varphi}(x)
  -
  \int \dd^Dx\,
  \delta\varphi(x)J_{\varphi}(x)        \\
  &\hspace{2.2cm}
  -
  \int \dd^Dx\,
  \varphi(x)\delta J_{\varphi}(x).
\end{aligned}
\]
The first and third terms cancel by the defining relation
\(\delta W/\delta J=\varphi\).  Therefore
\[
  \frac{\delta\Gamma[\varphi]}{\delta\varphi(x)}
  =
  -J_{\varphi}(x).
\]
At zero source, the physical vacuum expectation value is an extremum of the
effective action:
\[
  \left.
  \frac{\delta\Gamma[\varphi]}{\delta\varphi(x)}
  \right|_{\varphi=\vev{\phi}}
  =0 .
\]

The second derivative records the inverse relation between the connected
two-point function and the 1PI two-point kernel.  Define
\[
  W^{(2)}(x,y)
  =
  \frac{\delta^2 W[J]}{\delta J(x)\delta J(y)},
  \qquad
  \Gamma^{(2)}(x,y)
  =
  \frac{\delta^2\Gamma[\varphi]}{\delta\varphi(x)\delta\varphi(y)} .
\]
Differentiating \(\varphi(x)=\delta W/\delta J(x)\) with respect to
\(\varphi(y)\) gives
\[
  \delta^{(D)}(x-y)
  =
  \int \dd^Dz\,W^{(2)}(x,z)
  \frac{\delta J(z)}{\delta\varphi(y)} .
\]
Differentiating \(\delta\Gamma/\delta\varphi=-J\) gives
\[
  \Gamma^{(2)}(z,y)
  =
  -\frac{\delta J(z)}{\delta\varphi(y)} .
\]
Combining the two identities,
\[
  \int \dd^Dz\,
  W^{(2)}(x,z)\Gamma^{(2)}(z,y)
  =
  -\delta^{(D)}(x-y).
\]
The sign is a consequence of the Lorentzian convention
\(\Gamma=W-\int\varphi J\).  In momentum space, for a translation-invariant
vacuum,
\[
  W^{(2)}(k)\Gamma^{(2)}(k)=-1 .
\]
With the same Feynman \(+\ii0\) boundary prescription, \(W^{(2)}\) is related
to the connected propagator by the same factors of \(\ii\) that appear in the
source-derivative definition above.
\end{proof}

\section{Background Split and One-Particle-Irreducible Kernels}

\begin{proposition}[Shifted representation and tadpole condition]
\label{prop:shifted-1pi-representation-tadpole}
Fix a field configuration \(\varphi\), set \(J=J_{\varphi}\), and write
\[
  \phi=\varphi+\xi .
\]
The local 1PI functional has the shifted representation
\[
\begin{aligned}
  \ee^{\ii\Gamma[\varphi]}
  &=
  \ee^{\ii W[J]-\ii\langle J,\varphi\rangle}  \\
  &=
  \int [D\xi]\,
  \exp\!\left\{
    \ii S[\varphi+\xi]
    +
    \ii\langle J,\xi\rangle
  \right\}_{J=J_{\varphi}} .
\end{aligned}
\]
In this shifted functional integral,
\[
  \vev{\xi(x)}_{\varphi}=0 ,
\]
so the linear source term is chosen exactly to cancel the quantum tadpole in
the background \(\varphi\).
\end{proposition}

\begin{proof}
The first equality is the Legendre definition
\(\Gamma[\varphi]=W[J_\varphi]-\langle J_\varphi,\varphi\rangle\).  In the
regulated integral for \(Z[J_\varphi]\), make the affine change of variables
\(\phi=\varphi+\xi\).  The symbol \([D\xi]\) denotes the pushforward of the
specified regulated integration object under this change of variables; if the
chosen regulator produces a nontrivial Jacobian or density term, its logarithm
is part of the shifted action.  The source term
\(\langle J_\varphi,\phi\rangle\) becomes
\(\langle J_\varphi,\varphi\rangle+\langle J_\varphi,\xi\rangle\), and the
first part is precisely removed by the Legendre factor.  Finally,
\[
  \vev{\xi(x)}_{\varphi}
  =
  \vev{\phi(x)}_{J_\varphi}-\varphi(x)
  =
  \frac{\delta W[J]}{\delta J(x)}\bigg|_{J=J_\varphi}-\varphi(x)
  =
  0 .
\]
\end{proof}

Expand the action around the background:
\[
  S[\varphi+\xi]
  =
  S[\varphi]
  +
  \int \dd^Dx\,S^{(1)}_{\varphi}(x)\xi(x)
  +
  \frac12
  \int \dd^Dx\,\dd^Dy\,
  \xi(x)S^{(2)}_{\varphi}(x,y)\xi(y)
  +\cdots .
\]
The vertices in the shifted theory may carry two types of legs: quantum
\(\xi\)-legs that are contracted in loops, and background
\(\varphi\)-legs that remain external.  The source term supplies an additional
linear vertex
\[
  \ii\langle J_{\varphi},\xi\rangle .
\]
For example, if the microscopic action contains a cubic interaction, then
the polynomial identity after the split is
\[
  \phi^3
  =
  \varphi^3+3\varphi^2\xi
  +3\varphi\xi^2+\xi^3 .
\]
Thus a single local interaction produces vertices with zero, one, two, or
three quantum fluctuation legs.  The first term is a purely background
contribution to the zero-loop effective action; the terms with
\(\xi\)-legs are contracted in the shifted functional integral; and
the source supplies one more one-point quantum vertex.  The shifted-vertex
bookkeeping is displayed in Figure~\ref{fig:background-field-shifted-cubic-vertices}.
\begin{figure}[!ht]
\centering
  \begin{tikzpicture}[>=Stealth,scale=0.86,
    q/.style={line width=0.75pt},
    bg/.style={line width=0.75pt,densely dotted,blue!70!black},
    src/.style={line width=0.75pt,orange!85!black},
    dot/.style={circle,fill=black,inner sep=1.5pt}]
    \node[dot] (v) at (-4.4,0) {};
    \foreach \a in {25,145,270} \draw[q] (v) -- ++(\a:0.82);
    \node at (-4.4,-1.03) {\(\phi^3\)};
    \node at (-3.25,0) {\(\rightsquigarrow\)};
    \node[dot] (a) at (-2.1,0) {};
    \foreach \ang in {25,145,270} \draw[bg] (a) -- ++(\ang:0.75);
    \node at (-2.1,-1.03) {\(\varphi^3\)};
    \node[font=\large] at (-1.18,0) {\(+\)};
    \node[dot] (b) at (-0.2,0) {};
    \draw[q] (b) -- ++(270:0.78);
    \foreach \ang in {30,150} \draw[bg] (b) -- ++(\ang:0.78);
    \node at (-0.2,-1.03) {\(\varphi^2\xi\)};
    \node[font=\large] at (0.78,0) {\(+\)};
    \node[dot] (c) at (1.75,0) {};
    \draw[bg] (c) -- ++(90:0.78);
    \foreach \ang in {210,330} \draw[q] (c) -- ++(\ang:0.78);
    \node at (1.75,-1.03) {\(\varphi\xi^2\)};
    \node[font=\large] at (2.76,0) {\(+\)};
    \node[dot] (d) at (3.65,0) {};
    \foreach \ang in {25,145,270} \draw[q] (d) -- ++(\ang:0.78);
    \node at (3.65,-1.03) {\(\xi^3\)};
    \node[font=\large] at (4.7,0) {\(+\)};
    \draw[src] (5.25,0)--(5.95,0);
    \fill[orange!85!black] (5.95,0) circle (1.8pt);
    \node[orange!85!black] at (5.6,-0.48) {\(J_\varphi\xi\)};
  \end{tikzpicture}
\caption{Cubic interaction after the background-field split
\(\phi=\varphi+\xi\).  Dotted blue legs are background fields, solid legs are
quantum fluctuations, and the orange one-point vertex is the source term fixed
by the condition \(\langle\xi\rangle_\varphi=0\).}
\label{fig:background-field-shifted-cubic-vertices}
\end{figure}
Since \(J_{\varphi}\) is fixed by \(\vev{\xi}=0\), all tadpole
subdiagrams in the shifted theory cancel in the sum.

\begin{proposition}[One-particle-reducible graphs are Legendre trees]
\label{prop:opr-graphs-legendre-trees}
A connected diagram is one-particle reducible if cutting one internal
\(\xi\) propagator separates it into two connected diagrams, each
with at least one external background leg.  In the Legendre transform, such
single-propagator reducible contributions are generated by solving the
stationary equation for \(\varphi_J\) and joining lower kernels by exact
propagators.  They are therefore part of the tree reconstruction of connected
functions from \(\Gamma\), rather than new kernels in \(\Gamma\) itself.  The
new kernels in \(\Gamma\) are the one-particle-irreducible ones.
\end{proposition}

\begin{proof}
Equivalently, in a shifted connected vacuum diagram, cut an internal quantum
propagator.  Each side becomes a connected subdiagram with one open quantum
leg and any number of background legs.  When all graphs with the same
external background data are summed, either side is precisely the
source-dependent tadpole functional; it vanishes by
\(\vev{\xi}_\varphi=0\).  This is the diagrammatic content of the
Legendre transform, not an additional dynamical assumption.
\end{proof}

Figure~\ref{fig:volume-ii-tadpoles-and-1pi} shows the diagrammatic content of
this argument: cutting a quantum propagator in a shifted connected vacuum graph
produces tadpole subgraphs, which vanish at the Legendre-transform background.

\begin{figure}[!ht]
  \centering
  \resizebox{0.97\textwidth}{!}{%
  \begin{tikzpicture}[>=Stealth,x=1cm,y=1cm,every node/.style={font=\small}]
    \tikzset{
      qline/.style={line width=0.9pt},
      bg/.style={line width=1.0pt,blue!70!black,-{Latex[length=1.7mm]}},
      src/.style={line width=1.0pt,orange!80!black,-{Latex[length=1.7mm]}},
      blob/.style={circle,draw=black,fill=gray!17,line width=0.9pt,minimum size=0.85cm},
      smallblob/.style={circle,draw=black,fill=gray!17,line width=0.8pt,minimum size=0.58cm},
      arr/.style={-{Latex[length=2.0mm]},line width=0.9pt}
    }
    \begin{scope}
      \node[font=\normalsize] at (0,2.05) {tadpole condition};
      \node[smallblob] (T) at (0,0) {};
      \draw[qline] (T) .. controls (0.65,0.55) and (0.65,-0.55) .. (T);
      \draw[bg] (-1.65,0)--(T.west);
      \node[blue!70!black,left] at (-1.65,0) {\(\varphi\)};
      \node[font=\large] at (1.25,0) {\(+\)};
      \draw[src] (1.9,0)--(2.55,0);
      \node[orange!80!black,above] at (2.22,0.18) {\(J_{\varphi}\)};
      \node[font=\large] at (3.05,0) {\(=0\)};
      \node[align=center] at (1.2,-1.25) {\(\vev{\xi}_{\varphi}=0\)};
    \end{scope}

    \begin{scope}[shift={(7.2,0)}]
      \node[font=\normalsize] at (0,2.05) {one-particle reducible part};
      \node[blob] (A) at (-1.35,0) {};
      \node[blob] (B) at (1.35,0) {};
      \draw[qline] (A)--(B);
      \foreach \y in {-0.55,0.55} {
        \draw[bg] (-2.75,\y)--(A);
        \draw[bg] (2.75,\y)--(B);
      }
      \draw[red!70!black,dashed,line width=1pt] (0,-0.95)--(0,0.95);
      \node[red!70!black] at (0,-1.25) {single cut};
      \node[align=center] at (0,-1.85) {generated by trees\\of lower kernels};
    \end{scope}

    \begin{scope}[shift={(13.5,0)}]
      \node[font=\normalsize] at (0,2.05) {1PI kernel};
      \node[blob,minimum size=1.2cm] (C) at (0,0) {\(\Gamma^{(4)}\)};
      \foreach \ang in {35,145,215,325} {
        \draw[bg] (\ang:2.0)--(C);
      }
      \node[purple!80!black] at (0,-1.4) {irreducible block};
    \end{scope}
  \end{tikzpicture}
  }
  \caption{The source \(J_{\varphi}\) in the shifted functional integral is
  fixed by the condition that the quantum fluctuation has zero one-point
  function.  With this condition imposed, the Legendre transform is represented
  by one-particle-irreducible kernels.}
  \label{fig:volume-ii-tadpoles-and-1pi}
\end{figure}

Perturbatively, one may write
\[
  \ii\Gamma[\varphi]
  =
  \ii S[\varphi]
  +
  \left\{
    \hbox{sum of connected 1PI vacuum graphs in the shifted theory}
  \right\}.
\]
The position-space 1PI kernels are the functional derivatives
\[
  \Gamma^{(n)}(x_1,\ldots,x_n)
  =
  \left.
  \frac{\delta^n\Gamma[\varphi]}
       {\delta\varphi(x_1)\cdots\delta\varphi(x_n)}
  \right|_{\varphi=\varphi_0},
\]
where \(\varphi_0\) is the background about which the expansion is made.  In a
translation-invariant vacuum with vanishing one-point function, the expansion
point is \(\varphi_0=0\); if the one-point function is nonzero, one first
shifts the field and then expands around the shifted stationary background.
Equivalently, after expanding in powers of the
background, \(\ii\Gamma^{(n)}\) is computed by amputated 1PI \(n\)-point
diagrams whose external lines carry the background field.  The momenta on
these external legs range over the off-shell momentum space of the vertex
function.  Scattering kinematics enter only after LSZ and the corresponding
on-shell limits.

\section{Vertex Expansion}

\begin{definition}[Momentum-space 1PI vertices]
\label{def:momentum-space-1pi-vertices}
In a translation-invariant vacuum, the effective action has a momentum-space
expansion
\[
  \Gamma[\varphi]
  =
  \sum_{n=0}^{\infty}
  \frac1{n!}
  \int
  \prod_{r=1}^{n}
  \frac{\dd^Dk_r}{(2\pi)^D}\,
  (2\pi)^D\delta^{(D)}\!\left(\sum_{r=1}^{n}k_r\right)
  \Gamma^{(n)}(k_1,\ldots,k_n)\,
  \widetilde\varphi(k_1)\cdots\widetilde\varphi(k_n).
\]
The delta function expresses translation invariance, so only \(n-1\) of the
momenta are independent.  With this normalization,
\[
  \Gamma^{(n)}(k_1,\ldots,k_n)
  =
  \left.
  \frac{\delta^n\Gamma}
       {\delta\widetilde\varphi(k_1)\cdots
        \delta\widetilde\varphi(k_n)}
  \right|_{\varphi=0}
\]
after removing the overall momentum-conserving delta function.
\end{definition}

For \(n=2\), the exact two-point 1PI kernel is the inverse propagator in the
Legendre sense.  For a scalar field with tree-level action
\[
  S_0[\phi]
  =
  -\frac12
  \int \frac{\dd^Dk}{(2\pi)^D}\,
  \widetilde\phi(k)(k^2+m^2)\widetilde\phi(-k),
\]
one writes
\[
  \ii\Gamma^{(2)}(k)
  =
  -\ii(k^2+m^2)
  +
  \Pi_{\rm 1PI}(k),
\]
where \(\Pi_{\rm 1PI}(k)\) is the sum of 1PI self-energy diagrams in the same
normalization.  The exact connected propagator is obtained by inverting this
kernel with the \(i\epsilon\) prescription fixed by the underlying Lorentzian
functional integral.

\begin{definition}[Local derivative expansion of a 1PI branch]
\label{def:local-derivative-expansion-1pi-branch}
When the kernels are analytic near zero momentum, their Taylor expansion gives
a local derivative expansion.  For a single scalar field,
\[
  \Gamma[\varphi]
  =
  \int \dd^Dx\,
  \left[
    -V_{\rm eff}(\varphi)
    -\frac12 Z_{\rm eff}(\varphi)
      \partial_{\mu}\varphi\,\partial^{\mu}\varphi
    +\sum_{r\ge 2}
      \mathcal C_r(\varphi,\partial\varphi,\ldots)
  \right],
\]
where each \(\mathcal C_r\) contains operators with higher numbers of
derivatives.  The coefficients are determined by derivatives of the kernels
\(\Gamma^{(n)}(k_1,\ldots,k_n)\) at small momenta.  The expansion has to be
used with its analytic domain stated: loops of massless fields commonly
produce terms such as logarithms and threshold branch points at small
momentum, and these terms are part of the exact 1PI action.
\end{definition}

\section{One-Loop Effective Potential and Dimensional Transmutation}
\label{sec:coleman-weinberg-effective-potential}

The effective potential is the zero-derivative part of the local formal 1PI
effective action.  It is not, by definition, the exact convex
Legendre--Fenchel effective action of a positive Euclidean measure.  The latter
is discussed later in this chapter.  In perturbation theory one chooses a
background branch, computes the local 1PI functional around that branch, and
only then asks whether the branch is stable and whether the logarithms are
under perturbative control.

Work in four dimensions with one real scalar field and Euclidean action
\[
  S_E[\phi]
  =
  \int \dd^4x\,
  \left[
    \frac12\partial_\mu\phi\,\partial_\mu\phi
    +
    V_0(\phi)
  \right],
  \qquad
  V_0(\phi)=\frac{\lambda}{4!}\phi^4 .
\]
The regulated path-integral Lagrangian contains the counterterm part needed by
the chosen subtraction prescription; those terms are part of the regulated
Lagrangian, not a separate physical theory.  Let the background be constant,
\(\phi(x)=\varphi+\xi(x)\).  The quadratic fluctuation operator is
\[
  \Delta_\varphi
  =
  -\partial^2+M^2(\varphi),
  \qquad
  M^2(\varphi)=V_0''(\varphi)=\frac{\lambda}{2}\varphi^2 .
\]
In a finite box of volume \({\rm Vol}\), the formal one-loop contribution to
the 1PI effective action is
\[
  \Gamma_E^{(1)}[\varphi]
  =
  \frac12\operatorname{Tr}\log\Delta_\varphi ,
  \qquad
  V_1(\varphi)
  =
  \frac{1}{2\,{\rm Vol}}\operatorname{Tr}\log\Delta_\varphi .
\]
After passing to momentum space and subtracting the \(\varphi\)-independent
vacuum term,
\begin{equation}
  V_1(\varphi)
  =
  \frac12\mu^\epsilon
  \int\frac{\dd^{4-\epsilon}p}{(2\pi)^{4-\epsilon}}\,
  \log\bigl(p^2+M^2(\varphi)\bigr)
  +V_{\rm ct}^{(1)}(\varphi).
  \label{eq:cw-one-loop-log-determinant}
\end{equation}
The scale \(\mu\) is the dimensional-regularization scale in the selected
subtraction chart.  To evaluate the determinant without hiding constants,
differentiate with respect to \(M^2\):
\[
  \frac{\partial V_1}{\partial M^2}
  =
  \frac12\mu^\epsilon
  \int\frac{\dd^{4-\epsilon}p}{(2\pi)^{4-\epsilon}}\,
  \frac1{p^2+M^2}
  =
  \frac12
  \frac{\mu^\epsilon}{(4\pi)^{2-\epsilon/2}}
  \Gamma\!\left(-1+\frac{\epsilon}{2}\right)
  (M^2)^{1-\epsilon/2}.
\]
Integrating back in \(M^2\) and choosing the additive constant so that
\(V_1(0)=0\) gives
\[
  V_1^{\rm reg}
  =
  -\frac12
  \frac{\mu^\epsilon}{(4\pi)^{2-\epsilon/2}}
  \Gamma\!\left(-2+\frac{\epsilon}{2}\right)
  (M^2)^{2-\epsilon/2}
  +V_{\rm ct}^{(1)} .
\]
In the \(\overline{\rm MS}\) chart, after subtracting the pole and the
associated \(\log 4\pi-\gamma_E\) convention terms, this becomes
\begin{equation}
  V_1^{\overline{\rm MS}}(\varphi)
  =
  \frac{M^4(\varphi)}{64\pi^2}
  \left[
    \log\frac{M^2(\varphi)}{\mu^2}
    -\frac32
  \right].
  \label{eq:cw-one-loop-msbar-general}
\end{equation}
Thus the one-loop formal 1PI effective potential in massless \(\lambda\phi^4\)
theory is
\begin{equation}
  V_{\rm eff}^{\overline{\rm MS}}(\varphi)
  =
  \frac{\lambda(\mu)}{4!}\varphi^4
  +
  \frac{\lambda(\mu)^2\varphi^4}{256\pi^2}
  \left[
    \log\frac{\lambda(\mu)\varphi^2}{2\mu^2}
    -\frac32
  \right]
  +O(\lambda^3).
  \label{eq:cw-msbar-massless-phi4}
\end{equation}
The logarithm is the first explicit signal that the classical scale-invariant
polynomial has become a renormalized family depending on a subtraction scale.
The dependence on \(\mu\) is not physical by itself; it is compensated by the
running coupling.  To this order the field anomalous dimension starts only at
higher loop order, while
\[
  \beta_\lambda(\lambda)
  =
  \mu\frac{\dd\lambda}{\dd\mu}
  =
  \frac{3\lambda^2}{16\pi^2}+O(\lambda^3).
\]
Indeed,
\[
  \mu\frac{\partial}{\partial\mu}
  \left[
    \frac{\lambda^2\varphi^4}{256\pi^2}
    \log\frac{\lambda\varphi^2}{2\mu^2}
  \right]
  =
  -\frac{\lambda^2\varphi^4}{128\pi^2},
\]
whereas
\[
  \beta_\lambda\frac{\partial}{\partial\lambda}
  \left(\frac{\lambda\varphi^4}{24}\right)
  =
  \frac{\lambda^2\varphi^4}{128\pi^2}.
\]
Thus
\[
  \left(
    \mu\frac{\partial}{\partial\mu}
    +
    \beta_\lambda\frac{\partial}{\partial\lambda}
    -
    \gamma_\phi\varphi\frac{\partial}{\partial\varphi}
  \right)
  V_{\rm eff}=0
\]
holds through order \(\lambda^2\), with
\(\gamma_\phi=O(\lambda^2)\) contributing only at the next order in this
single-scalar example.

Coleman--Weinberg often uses a finite subtraction condition at a nonzero field
value \(M\), because the fourth derivative of \(\varphi^4\log\varphi^2\) is
singular at \(\varphi=0\).  Define \(\lambda_{\rm CW}\) by
\[
  \left.
  \frac{\dd^4 V_{\rm eff}}{\dd\varphi^4}
  \right|_{\varphi=M}
  =
  \lambda_{\rm CW}.
\]
After the corresponding finite change of quartic coordinate, the one-loop
potential is
\begin{equation}
  V_{\rm CW}(\varphi)
  =
  \frac{\lambda_{\rm CW}}{24}\varphi^4
  +
  \frac{\lambda_{\rm CW}^2}{256\pi^2}\varphi^4
  \left[
    \log\frac{\varphi^2}{M^2}
    -\frac{25}{6}
  \right]
  +O(\lambda_{\rm CW}^3).
  \label{eq:cw-renormalization-condition-potential}
\end{equation}
The constant \(-25/6\) is fixed by the displayed fourth-derivative condition:
\[
  \left.
  \frac{\dd^4}{\dd\varphi^4}
  \left[
    \varphi^4\log\frac{\varphi^2}{M^2}
  \right]
  \right|_{\varphi=M}
  =
  100,
  \qquad
  \frac{\dd^4}{\dd\varphi^4}\varphi^4=24 .
\]

The nonzero stationary points of the one-loop branch satisfy
\begin{equation}
  0
  =
  \frac{1}{\varphi^3}\frac{\dd V_{\rm CW}}{\dd\varphi}
  =
  \frac{\lambda_{\rm CW}}{6}
  +
  \frac{\lambda_{\rm CW}^2}{64\pi^2}
  \left[
    \log\frac{\varphi^2}{M^2}
    -\frac{11}{3}
  \right]
  +O(\lambda_{\rm CW}^3).
  \label{eq:cw-stationary-condition}
\end{equation}
At the one-loop level this gives the formal scale
\begin{equation}
  v
  =
  M
  \exp\!\left[
    \frac{11}{6}
    -
    \frac{16\pi^2}{3\lambda_{\rm CW}(M)}
  \right].
  \label{eq:cw-formal-generated-scale}
\end{equation}
At such a stationary point,
\[
  \left.
  \frac{\dd^2V_{\rm CW}}{\dd\varphi^2}
  \right|_{\varphi=v}
  =
  \frac{\lambda_{\rm CW}^2}{32\pi^2}v^2+O(\lambda_{\rm CW}^3v^2).
\]
Equations \eqref{eq:cw-stationary-condition} and
\eqref{eq:cw-formal-generated-scale} display the algebraic mechanism called
dimensional transmutation: a dimensionless coupling specified at the
subtraction point is exchanged for a dimensionful scale on a chosen
renormalized trajectory.

This statement has an important perturbative qualification.  In the single
real scalar theory, taking \(M=v\) in
\eqref{eq:cw-stationary-condition} gives
\[
  \lambda_{\rm CW}(v)=\frac{32\pi^2}{11},
\]
which is not a small coupling.  If instead \(\lambda_{\rm CW}(M)\) is small,
then \(\log(v/M)=O(1/\lambda_{\rm CW})\), so higher leading logarithms are not
parametrically smaller than the one-loop logarithm.  The one-loop
massless-scalar potential therefore demonstrates the structure of the
calculation and the scale-generating relation, but it does not by itself give
a controlled weak-coupling proof of a radiatively generated vacuum in pure
\(\lambda\phi^4\) theory.

The renormalization-group form of dimensional transmutation is independent of
this particular one-loop branch.  For a dimensionless coupling \(g\) with
\(\beta(g)\neq0\), choosing \(g(\mu)\) at a subtraction scale is equivalent,
on a one-dimensional RG trajectory, to choosing the scale \(\Lambda_{g_*}\)
at which the running coupling reaches a reference value \(g_*\):
\begin{equation}
  \log\frac{\Lambda_{g_*}}{\mu}
  =
  \int_{g(\mu)}^{g_*}\frac{\dd g}{\beta(g)}.
  \label{eq:dimensional-transmutation-rg-scale}
\end{equation}
For the one-loop scalar beta function \( \beta_\lambda=b\lambda^2\),
\(b=3/(16\pi^2)\), this gives
\[
  \Lambda_{\lambda_*}
  =
  \mu
  \exp\!\left[
    \frac{1}{b\lambda(\mu)}
    -
    \frac{1}{b\lambda_*}
  \right].
\]
Because \(b>0\), this scale is the scale of increasing coupling toward the
ultraviolet rather than an asymptotically free strong scale.  In gauge theories
with a negative one-loop beta function the same formula produces the familiar
infrared strong scale.

A controlled Coleman--Weinberg minimum can occur in a multi-coupling theory
where the tree quartic coupling is of the same order as a one-loop effect.
For example, in massless scalar electrodynamics on a chosen perturbative
branch and in the conventional radial-field normalization, the gauge-boson
one-loop term gives
\[
  V_{\rm eff}(\varphi)
  =
  \frac{\lambda}{24}\varphi^4
  +
  \frac{3e^4}{64\pi^2}\varphi^4
  \left[
    \log\frac{\varphi^2}{M^2}
    -\frac{25}{6}
  \right]
  +\cdots ,
\]
after the same type of nonzero-field quartic subtraction.  Then a stationary
point at \(M=v\) requires \(\lambda=33e^4/(8\pi^2)+O(e^6)\), which is
compatible with weak coupling when \(e\) is small.  This is the controlled
version of radiative breaking; the single-scalar computation above is the
minimal determinant calculation and the clearest place to see how the
subtraction scale enters the formal 1PI potential.

\section{Connected Functions from Exact 1PI Vertices}

The Legendre transform can be inverted.  Given \(\Gamma[\varphi]\), the
connected functional \(W[J]\) is obtained by solving
\[
  \frac{\delta\Gamma[\varphi]}{\delta\varphi(x)}
  =
  -J(x)
\]
for \(\varphi=\varphi_J\), and then setting
\[
  W[J]
  =
  \Gamma[\varphi_J]
  +
  \langle J,\varphi_J\rangle .
\]
Equivalently, the inversion can be written as the stationary value
\[
  W[J]
  =
  \operatorname*{stationary}_{\varphi}
  \left(
    \Gamma[\varphi]
    +
    \langle J,\varphi\rangle
  \right).
\]
This identity is exact at the level of the Legendre transform.  It also has a
diagrammatic interpretation: connected Green functions are obtained as tree
diagrams whose vertices are exact 1PI kernels and whose internal lines are the
exact connected propagator.

\begin{proposition}[Exact 1PI tree reconstruction of connected kernels]
\label{prop:exact-1pi-tree-reconstruction}
Work in DeWitt notation, with repeated indices summed and integrated, and set
\[
  W^{i_1\cdots i_n}
  =
  \frac{\delta^nW}{\delta J_{i_1}\cdots\delta J_{i_n}},
  \qquad
  \Gamma_{i_1\cdots i_n}
  =
  \frac{\delta^n\Gamma}
       {\delta\varphi^{i_1}\cdots\delta\varphi^{i_n}}
\]
at the paired configurations \(J\) and \(\varphi_J\).  Let
\[
  G^{ij}=W^{ij}
  =
  \frac{\delta\varphi_J^i}{\delta J_j}
\]
be the exact connected two-point kernel.  Then
\begin{equation}
  \Gamma_{ia}G^{aj}=-\delta_i{}^j,
  \label{eq:exact-1pi-tree-two-point}
\end{equation}
\begin{equation}
  W^{ijk}
  =
  G^{ia}G^{jb}G^{kc}\Gamma_{abc},
  \label{eq:exact-1pi-tree-three-point}
\end{equation}
and
\begin{equation}
\begin{aligned}
  W^{ijkl}
  &=
  G^{ia}G^{jb}G^{kc}G^{ld}\Gamma_{abcd}
\\
  &\quad
  +G^{ia}G^{jb}\Gamma_{abe}G^{ef}\Gamma_{fcd}G^{kc}G^{ld}
\\
  &\quad
  +G^{ia}G^{kc}\Gamma_{ace}G^{ef}\Gamma_{fbd}G^{jb}G^{ld}
\\
  &\quad
  +G^{ia}G^{ld}\Gamma_{ade}G^{ef}\Gamma_{fbc}G^{jb}G^{kc}.
\end{aligned}
  \label{eq:exact-1pi-tree-four-point}
\end{equation}
In general, every connected kernel \(W^{i_1\cdots i_n}\) is obtained by a
finite sum over trees whose vertices are exact 1PI kernels
\(\Gamma_r\) with \(r\ge3\), whose internal edges are \(G\), and whose
external half-edges are attached to the external labels by \(G\).  No loop is
created in this reconstruction; all loop information is already contained in
the exact kernels \(\Gamma_r\) and the exact propagator \(G\).
\end{proposition}

\begin{proof}
The stationarity equation is
\[
  \Gamma_i[\varphi_J]+J_i=0.
\]
Differentiating once with respect to \(J_j\) gives
\[
  \Gamma_{ia}G^{aj}+\delta_i{}^j=0,
\]
which proves \eqref{eq:exact-1pi-tree-two-point}.  Differentiating this
identity with respect to \(J_k\) gives
\[
  \Gamma_{iab}G^{bk}G^{aj}
  +
  \Gamma_{ia}W^{ajk}
  =
  0 .
\]
Multiplying by \(G^{\ell i}\) and using
\(G^{\ell i}\Gamma_{ia}=-\delta^\ell{}_a\), which is the same inverse
relation with the two factors interchanged, gives
\[
  W^{\ell jk}
  =
  G^{\ell i}\Gamma_{iab}G^{aj}G^{bk}.
\]
Renaming dummy indices proves
\eqref{eq:exact-1pi-tree-three-point}.

Differentiate the preceding second-derivative identity once more:
\[
\begin{aligned}
  0
  &=
  \Gamma_{iabc}G^{c\ell}G^{bk}G^{aj}
  +
  \Gamma_{iab}W^{bk\ell}G^{aj}
  +
  \Gamma_{iab}G^{bk}W^{aj\ell}
\\
  &\quad
  +
  \Gamma_{iab}G^{b\ell}W^{ajk}
  +
  \Gamma_{ia}W^{ajk\ell}.
\end{aligned}
\]
Multiplication by \(G^{mi}\) solves for \(W^{mjk\ell}\).  The first term gives
the single four-point vertex in
\eqref{eq:exact-1pi-tree-four-point}; substituting
\eqref{eq:exact-1pi-tree-three-point} into the other three terms gives the
three pairings of two exact three-point vertices joined by one exact
propagator.  This proves \eqref{eq:exact-1pi-tree-four-point}.

The recursive tree statement follows by induction.  Suppose every source
derivative of \(\varphi_J\) of order less than \(n\) is a sum of trees.  The
\((n-1)\)st derivative of
\(\Gamma_i[\varphi_J]+J_i=0\) is a finite Fa\`a di Bruno sum: one differentiates
some \(\Gamma_r\) and distributes the source derivatives among the factors
\(\delta^{m}\varphi_J/\delta J^m\).  By the induction hypothesis those
factors are trees.  Multiplying once by \(G\) solves the remaining
\(\Gamma_{ia}\)-factor and attaches one additional external propagator.  This
operation glues trees at a new exact 1PI vertex and cannot form a cycle.
Therefore the result is a sum over trees only.  Translating from derivatives
of \(W\) to Lorentzian Green functions uses the same powers of \(\ii\) as in
the source-derivative definitions at the beginning of the chapter.
\end{proof}

\begin{figure}[!ht]
  \centering
  \resizebox{0.95\textwidth}{!}{%
  \begin{tikzpicture}[>=Stealth,x=1cm,y=1cm,every node/.style={font=\small}]
    \tikzset{
      leg/.style={line width=0.95pt},
      blob/.style={circle,draw=black,fill=purple!12,line width=0.9pt,minimum size=0.78cm},
      arr/.style={-{Latex[length=2.0mm]},line width=0.9pt}
    }
    \node[font=\normalsize] at (0,1.9) {connected four-point function};
    \node[blob,fill=gray!16,minimum size=1.15cm] (C) at (0,0) {};
    \foreach \ang in {35,145,215,325} {
      \draw[leg] (\ang:2.0)--(C);
    }
    \node at (0,-1.65) {\(G^{(4)}_{\rm c}\)};

    \draw[arr] (2.35,0) -- (4.1,0);

    \begin{scope}[shift={(6.2,0)}]
      \node[blob] (V4) at (-1.6,0) {\(\Gamma^{(4)}\)};
      \foreach \ang in {35,145,215,325} {
        \draw[leg] ($ (V4) + (\ang:1.55) $)--(V4);
      }
      \node[font=\large] at (0.1,0) {\(+\)};
      \node[blob] (V3a) at (1.25,0.55) {\(\Gamma^{(3)}\)};
      \node[blob] (V3b) at (3.05,-0.55) {\(\Gamma^{(3)}\)};
      \draw[leg] (V3a)--(V3b);
      \draw[leg] ($(V3a)+(-0.95,0.72)$)--(V3a);
      \draw[leg] ($(V3a)+(-0.95,-0.72)$)--(V3a);
      \draw[leg] ($(V3b)+(0.95,0.72)$)--(V3b);
      \draw[leg] ($(V3b)+(0.95,-0.72)$)--(V3b);
      \node[font=\large] at (4.55,0) {\(+\cdots\)};
      \node at (1.6,-1.65) {trees only; vertices and lines are exact};
    \end{scope}
  \end{tikzpicture}
  }
  \caption{Connected Green functions are reconstructed from the exact 1PI
  action by forming trees.  The internal lines are exact propagators, and the
  vertices are exact kernels \(\Gamma^{(n)}\).}
  \label{fig:volume-ii-connected-from-1pi}
\end{figure}

One can express the same statement with a bookkeeping parameter \(\hbar\):
\[
  \ii W[J]
  =
  \lim_{\hbar\to0}
  \hbar\log
  \int [D\varphi]\,
  \exp\!\left\{
    \frac{\ii}{\hbar}
    \left(
      \Gamma[\varphi]
      +
      \langle J,\varphi\rangle
    \right)
  \right\}.
\]
The limit keeps only tree diagrams built from \(\Gamma\).  Since
\(\Gamma\) already contains the full loop information of the original theory,
this tree expansion is the tree-level representation of the exact Legendre
inverse.

\section{Schwinger--Dyson Identities and the 1PI Hierarchy}
\label{sec:schwinger-dyson-identities-1pi-hierarchy}

The Legendre transform organizes exact kernels, but it does not by itself use
the equation of motion of the integration variable.  The latter information is
contained in the Schwinger--Dyson identities.  These identities are exact only
after their hypotheses are stated: at the regulator under discussion the
reference integration density must be translation invariant in the integration
coordinate, or else its logarithmic derivative must be included; the contour or
domain must have no boundary contribution for the variation being made; and
any Jacobian from a nontrivial change of variables must be recorded as an
additional insertion.  Gauge theories replace the naive field translation by
gauge-fixed BRST/BV identities.  In this scalar section we assume the
translation-invariant regulated scalar density introduced at the beginning of
the chapter.

Use condensed DeWitt notation in this section.  A single index \(i\) denotes
both a discrete field label and a spacetime point, repeated indices are summed
and integrated, and
\[
  J_i\phi^i=\langle J,\phi\rangle .
\]
Write
\[
  S_i[\phi]=\frac{\delta S[\phi]}{\delta\phi^i},
  \qquad
  S_{ij}[\phi]=\frac{\delta^2S[\phi]}
    {\delta\phi^i\delta\phi^j},
\]
and similarly for higher field derivatives of the regulated action.

\begin{proposition}[Functional Schwinger--Dyson identity]
\label{prop:functional-schwinger-dyson-identity}
For every regulated test functional \(\mathcal O[\phi]\) for which the
following integration by parts is legitimate,
\begin{equation}
  \left\langle
    \bigl(S_i[\phi]+J_i\bigr)\mathcal O[\phi]
  \right\rangle_J
  =
  \ii\left\langle
    \frac{\delta\mathcal O[\phi]}{\delta\phi^i}
  \right\rangle_J .
  \label{eq:sd-insertion-identity}
\end{equation}
Equivalently, at the level of the source functional,
\begin{equation}
  \left[
    S_i\!\left(\frac1{\ii}\frac{\delta}{\delta J}\right)
    +J_i
  \right]Z[J]=0 .
  \label{eq:sd-functional-z}
\end{equation}
\end{proposition}

\begin{proof}
At finite regulator, integrate the total derivative
\[
  0=
  \int D_\Lambda^{\rm ref}\phi\,
  \frac{\delta}{\delta\phi^i}
  \left[
    \mathcal O[\phi]\,
    \exp\{\ii S[\phi]+\ii J_j\phi^j\}
  \right].
\]
Expanding the derivative gives
\[
  0=
  \left\langle
    \frac{\delta\mathcal O}{\delta\phi^i}
  \right\rangle_J
  +
  \ii
  \left\langle
    \mathcal O(S_i+J_i)
  \right\rangle_J .
\]
Since \(-1/\ii=\ii\), this is
\eqref{eq:sd-insertion-identity}.  Taking \(\mathcal O=1\) and replacing each
field insertion by \((1/\ii)\delta/\delta J\) gives
\eqref{eq:sd-functional-z}.  If the reference density has logarithmic
derivative \(R_i[\phi]\), the left side is modified by
\(S_i+J_i-\ii R_i\); that is the measure contribution which, in symmetry
variations, becomes the anomaly term.
\end{proof}

The case \(\mathcal O=1\) gives the quantum equation of motion
\[
  \left\langle S_i[\phi]\right\rangle_J+J_i=0.
\]
Using \(\Gamma_i[\varphi]=\delta\Gamma/\delta\varphi^i=-J_i\), this becomes
\begin{equation}
  \Gamma_i[\varphi]
  =
  \left\langle S_i[\phi]\right\rangle_{J_\varphi}.
  \label{eq:sd-1pi-equation-of-motion}
\end{equation}
Thus the 1PI equation of motion is not the classical equation
\(S_i[\varphi]=0\).  It is the expectation value of the microscopic equation
of motion in the source ensemble whose mean field is \(\varphi\).

Equation \eqref{eq:sd-functional-z} may be written directly in connected
variables.  Since
\[
  Z^{-1}\frac1{\ii}\frac{\delta Z}{\delta J_i}
  =
  \frac{\delta W}{\delta J_i}
  =
  \varphi^i,
\]
one has the operator identity
\begin{equation}
  Z[J]^{-1}
  F\!\left(\frac1{\ii}\frac{\delta}{\delta J}\right)Z[J]
  =
  F\!\left(\varphi+\frac1{\ii}\frac{\delta}{\delta J}\right)1
  \label{eq:source-derivative-shift-identity}
\end{equation}
for every polynomial functional \(F\), with the functional derivatives on the
right acting on all source-dependent quantities to their right before setting
them on \(1\).  Therefore
\begin{equation}
  S_i\!\left(\varphi+\frac1{\ii}\frac{\delta}{\delta J}\right)1
  +J_i=0 .
  \label{eq:sd-connected-functional-form}
\end{equation}
This formula is the compact exact Schwinger--Dyson equation for \(W\).  Its
Legendre form is obtained by substituting \(J_i=-\Gamma_i[\varphi]\) and
using \(W^{(2)}\Gamma^{(2)}=-1\) to convert source derivatives into field
derivatives when desired.

\subsection*{Quartic scalar example}

For a single scalar field with a quartic interaction, write the regulated
action in condensed notation as
\[
  S[\phi]
  =
  \frac12 K_{ij}\phi^i\phi^j
  +
  \frac{\lambda}{4!}V_{ijkl}\phi^i\phi^j\phi^k\phi^l ,
\]
where \(K_{ij}=K_{ji}\) is the regulated quadratic kernel and
\(V_{ijkl}\) is totally symmetric.  The sign of \(\lambda V\) is the sign
chosen in the Lorentzian action; no Euclidean positivity is being assumed in
this formal identity.  Since
\[
  S_i[\phi]=K_{ij}\phi^j+\frac{\lambda}{3!}V_{ijkl}\phi^j\phi^k\phi^l,
\]
the connected Schwinger--Dyson equation becomes
\begin{equation}
\begin{aligned}
  \Gamma_i[\varphi]
  &=
  K_{ij}\varphi^j
  +
  \frac{\lambda}{3!}V_{ijkl}
  \Bigl[
    \varphi^j\varphi^k\varphi^l   \\
  &\qquad\qquad
    +\frac1{\ii}
      \bigl(
        W^{jk}\varphi^l
        +W^{jl}\varphi^k
        +W^{kl}\varphi^j
      \bigr)
    +\frac1{\ii^2}W^{jkl}
  \Bigr]_{J=J_\varphi}.
\end{aligned}
  \label{eq:sd-quartic-1pi-field-equation}
\end{equation}
Here
\[
  W^{i_1\cdots i_n}
  =
  \frac{\delta^nW[J]}
       {\delta J_{i_1}\cdots\delta J_{i_n}}
\]
is the connected-kernel convention already used in this chapter.  The powers
of \(1/\ii\) in \eqref{eq:sd-quartic-1pi-field-equation} are not optional:
they come from the Lorentzian source factor \(\exp(\ii J_i\phi^i)\).

\subsection*{Full propagator member}

Assume for this paragraph that the vacuum is translation invariant and
\(\phi\mapsto-\phi\) symmetric, so that
\(\varphi=0\) and odd connected functions vanish at \(J=0\).  Differentiate
the field equation with respect to \(J_m\) and set \(J=0\).  The exact
two-point member of the hierarchy is
\begin{equation}
\begin{aligned}
  0
  &=
  K_{ij}W^{jm}
  +
  \frac{\lambda}{3!}V_{ijkl}
  \Bigl[
    \frac1{\ii}
      \bigl(
        W^{jk}W^{lm}
        +W^{jl}W^{km}
        +W^{jm}W^{kl}
      \bigr)
    +\frac1{\ii^2}W^{jklm}
  \Bigr]
  +\delta_i{}^m .
\end{aligned}
  \label{eq:sd-full-propagator-connected}
\end{equation}
This is the full propagator Schwinger--Dyson equation in connected form.  It
is not a one-loop formula: \(W^{jm}\) is the exact connected two-point kernel
and \(W^{jklm}\) is the exact connected four-point kernel.  Multiplying
\eqref{eq:sd-full-propagator-connected} by the exact inverse kernel
\(\Gamma^{(2)}\), using
\[
  W^{im}\Gamma^{(2)}_{mn}=-\delta^i{}_n,
\]
gives the corresponding inverse-propagator form
\begin{equation}
  \Gamma^{(2)}_{in}
  =
  K_{in}
  +
  \frac{\lambda}{2\ii}V_{inkl}W^{kl}
  -
  \frac{\lambda}{3!\,\ii^2}
  V_{ijkl}W^{jklm}\Gamma^{(2)}_{mn}.
  \label{eq:sd-full-propagator-inverse}
\end{equation}
The second term is the exact tadpole contraction.  The last term is the
amputated contribution of the exact connected four-point kernel.  A common
``gap'' or Hartree approximation drops the connected four-point term and then
solves the remaining equation self-consistently for \(W^{(2)}\).  That
operation is a truncation of the exact hierarchy, not a derivation from the
Schwinger--Dyson identity.

\subsection*{Full vertex member}

The next nontrivial member is obtained from
\eqref{eq:sd-insertion-identity} with
\(\mathcal O=\phi^a\phi^b\phi^c\), or equivalently by differentiating
\eqref{eq:sd-quartic-1pi-field-equation} three times with respect to the
source.  First, the insertion form gives
\begin{equation}
\begin{aligned}
  &K_{ij}
    \left\langle\phi^j\phi^a\phi^b\phi^c\right\rangle
  +
  \frac{\lambda}{3!}V_{ijkl}
    \left\langle
      \phi^j\phi^k\phi^l\phi^a\phi^b\phi^c
    \right\rangle     \\
  &\qquad =
  \ii\bigl(
    \delta_i{}^a\left\langle\phi^b\phi^c\right\rangle
    +\delta_i{}^b\left\langle\phi^a\phi^c\right\rangle
    +\delta_i{}^c\left\langle\phi^a\phi^b\right\rangle
  \bigr).
\end{aligned}
  \label{eq:sd-full-four-point-member}
\end{equation}
This equation contains disconnected pair products.  They are not an ambiguity:
they are removed by the lower two-point Schwinger--Dyson member.  The same
statement is obtained without mentioning disconnected moments by differentiating
the connected equation.  Let
\(\operatorname{Sym}_{abc}\) denote summation over all distinct permutations
of the displayed labels \(a,b,c\), and let
\(\operatorname{Sym}^{2,1}_{abc}\) denote the three splittings
\((ab|c),(ac|b),(bc|a)\).  In the symmetric vacuum,
\begin{equation}
\begin{aligned}
  0
  &=
  K_{ij}W^{jabc}
  +
  \frac{\lambda}{3!}V_{ijkl}
  \Bigl[
    \operatorname{Sym}_{abc}
      W^{ja}W^{kb}W^{lc}                                      \\
  &\qquad\qquad
    +\frac{3}{\ii}
      \Bigl(
        W^{jk}W^{labc}
        +
        \operatorname{Sym}^{2,1}_{abc}
        W^{jkab}W^{lc}
      \Bigr)
    +\frac1{\ii^2}W^{jklabc}
  \Bigr] .
\end{aligned}
  \label{eq:sd-connected-four-point-member}
\end{equation}
Equation \eqref{eq:sd-connected-four-point-member} is often the cleanest form
of the full vertex hierarchy: every term is connected, and the first term is
the kinetic operator acting on the exact connected four-point kernel.  The
1PI version is obtained by substituting the Legendre identities between
connected kernels and exact 1PI kernels.  In the
\(\phi\mapsto-\phi\) symmetric case,
\begin{equation}
  W^{abcd}
  =
  W^{aa'}W^{bb'}W^{cc'}W^{dd'}
  \Gamma^{(4)}_{a'b'c'd'} ,
  \label{eq:sd-four-point-amputation-z2}
\end{equation}
because the exact three-point vertex vanishes.  The connected six-point
kernel decomposes as
\begin{equation}
\begin{aligned}
  W^{abcdef}
  &=
  W^{aa'}W^{bb'}W^{cc'}W^{dd'}W^{ee'}W^{ff'}
  \Gamma^{(6)}_{a'b'c'd'e'f'}                                      \\
  &\quad
  +
  \sum_{\mathcal P_3}
  \left(
    \prod_{r\in A}W^{rr'}
  \right)
  \left(
    \prod_{s\in B}W^{ss'}
  \right)
  \Gamma^{(4)}_{A'u}
  W^{uv}
  \Gamma^{(4)}_{vB'} .
\end{aligned}
  \label{eq:sd-six-point-legendre-z2}
\end{equation}
Here \(\mathcal P_3\) is the set of the ten unordered partitions
\(\{a,b,c,d,e,f\}=A\sqcup B\) with \(|A|=|B|=3\); \(A'\) and \(B'\) denote
the primed external indices attached to the labels in the corresponding
block.  The signs in
\eqref{eq:sd-four-point-amputation-z2} and
\eqref{eq:sd-six-point-legendre-z2} follow from the convention
\(W^{im}\Gamma^{(2)}_{mn}=-\delta^i{}_n\).  Substituting
\eqref{eq:sd-four-point-amputation-z2} and
\eqref{eq:sd-six-point-legendre-z2} into
\eqref{eq:sd-connected-four-point-member}, and multiplying by exact inverse
two-point kernels on the external labels, gives the full four-point 1PI
Schwinger--Dyson equation.  It relates \(\Gamma^{(4)}\) to the exact
propagator, the exact six-point vertex \(\Gamma^{(6)}\), and the two
\(\Gamma^{(4)}\)-vertices joined by one exact propagator in
\eqref{eq:sd-six-point-legendre-z2}.  Without the
\(\phi\mapsto-\phi\) symmetry, the same construction contains the additional
tree terms built from exact three-point vertices described in
Figure~\ref{fig:volume-ii-connected-from-1pi}.  A vertex truncation is the
choice to keep only a specified subset of these terms, for example to model
\(\Gamma^{(4)}\) while discarding the connected six-point kernel.  The exact
Schwinger--Dyson hierarchy itself does not justify that closure; the closure
must be supplied by a small parameter, a controlled kinematic limit, a
renormalization condition plus error estimate, or an explicit nonperturbative
ansatz.

\section{Euclidean Convexity}

The convexity statement is a theorem only after the Euclidean functional
integral has been given positive measure-theoretic meaning.  The analytic input
used in the proof is isolated in the following assumption.

\begin{assumption}[Positive finite Euclidean regulator for convexity]
\label{ass:positive-euclidean-regulator-convexity}
Fix a regulator \(\Lambda\).  There is a measurable regulated configuration
space \((\mathcal C_\Lambda,\mathcal A_\Lambda)\), a positive finite measure
\(\mu_\Lambda\) on it, and a real vector space \(\mathcal S_\Lambda\) of
regulated sources paired with the regulated field by a measurable real
functional
\[
  (J,\phi)\longmapsto \langle J,\phi\rangle_\Lambda .
\]
There is a nonempty convex source domain
\(\mathcal D_\Lambda\subset\mathcal S_\Lambda\) such that, for every
\(J\in\mathcal D_\Lambda\),
\begin{equation}
  0<
  Z_{E,\Lambda}[J]
  :=
  \int_{\mathcal C_\Lambda}
  \exp\{-\langle J,\phi\rangle_\Lambda\}\,
  \dd\mu_\Lambda(\phi)
  <\infty .
  \label{eq:positive-regulator-partition-function}
\end{equation}
In the scalar finite-dimensional convention this positive measure may be
represented as a density
\[
  \dd\mu_\Lambda(\phi)
  =
  Z_\Lambda^{-1}
  \ee^{-S_E[\phi]}\,
  D_\Lambda^{\rm ref}\phi ,
\]
where the normalization is optional in unnormalized source functionals.  In the
Gaussian-reference convention the same object is written
\(\ee^{-L_\Lambda[\phi]}\dd\mu_{C_\Lambda}(\phi)\), with the relation to
\(D_\Lambda^{\rm ref}\phi\) given by
\eqref{eq:gaussian-reference-measure-from-density}.
\end{assumption}

Under Assumption~\ref{ass:positive-euclidean-regulator-convexity} define
\[
  Z_E[J]=Z_{E,\Lambda}[J]=\ee^{-W_E[J]} .
\]
Any continuum limit \(\Lambda\to\infty\) of the statements below requires an
additional existence theorem for the limiting Schwinger functional.  The
finite-regulator argument by itself does not prove such a limit for a proposed
model.

\begin{definition}[Lorentzian--Euclidean source-convention bridge]
\label{def:lorentzian-euclidean-source-bridge}
Write the Lorentzian logarithm as
\[
  \mathcal W_L[J_L]=\log Z_L[J_L]=\ii W_L[J_L],
  \qquad
  Z_L[J_L]
  =
  \int[D\phi]\,\exp\!\left\{\ii S_L[\phi]
  +\ii\langle J_L,\phi\rangle_L\right\}.
\]
Write the Euclidean logarithm in the convention used in this and the following
Euclidean chapters as
\[
  \mathcal W_E[J_E]=\log Z_E[J_E]=-W_E[J_E],
  \qquad
  Z_E[J_E]
  =
  \int[D\phi]\,\exp\!\left\{-S_E[\phi]-\langle J_E,\phi\rangle_E\right\}.
\]
In this definition \([D\phi]\) is only compact continuum notation for the
regulated integration object already specified earlier in the chapter.  In a
finite-dimensional bosonic scalar regulator, the Euclidean line is represented
by \(D_\Lambda^{\rm ref}\phi\,\ee^{-S_{E,\Lambda}[\phi]}\); after splitting
\(S_{E,\Lambda}=S_{\rm kin,\Lambda}+L_\Lambda\), the same object is
\(\dd\mu_{C_\Lambda}(\phi)\,\ee^{-L_\Lambda[\phi]}\).  The Lorentzian line is
the corresponding oscillatory notation and carries its own regulator and
boundary-value prescription.
For scalar insertions with no tensor time indices, the analytic continuation
from Lorentzian time-ordered connected distributions to Euclidean connected
Schwinger distributions is recorded by the derivative identities
\begin{equation}
  G^{L,{\rm c}}_n
  =
  \left.
  \ii^{-n}
  \frac{\delta^n\mathcal W_L}
       {\delta J_L(x_1)\cdots\delta J_L(x_n)}
  \right|_{J_L=0}
  =
  \left.
  \ii^{1-n}
  \frac{\delta^n W_L}
       {\delta J_L(x_1)\cdots\delta J_L(x_n)}
  \right|_{J_L=0},
  \label{eq:lorentzian-connected-derivative-normalization}
\end{equation}
and
\begin{equation}
  G^{E,{\rm c}}_n
  =
  \left.
  (-1)^n
  \frac{\delta^n\mathcal W_E}
       {\delta J_E(x_1)\cdots\delta J_E(x_n)}
  \right|_{J_E=0}
  =
  \left.
  (-1)^{n+1}
  \frac{\delta^n W_E}
       {\delta J_E(x_1)\cdots\delta J_E(x_n)}
  \right|_{J_E=0}.
  \label{eq:euclidean-connected-derivative-normalization}
\end{equation}
The source signs in the exponents are related as follows.  Under the formal
coordinate rotation \(t=-\ii\tau\), the Lorentzian pairing satisfies
\[
  \ii\langle J_L,\phi\rangle_L
  \longmapsto
  +\langle J_L^{\rm W},\phi_E\rangle_E,
  \qquad
  J_L^{\rm W}(\tau,\vec x)=J_L(-\ii\tau,\vec x),
\]
whereas the Euclidean convention above uses
\(-\langle J_E,\phi_E\rangle_E\).  Therefore the Euclidean source variable of
this chapter is
\begin{equation}
  J_E=-J_L^{\rm W}
  \label{eq:lorentzian-euclidean-source-sign-map}
\end{equation}
when one translates a formal source insertion from the Lorentzian convention
to the Euclidean convention.  Equations
\eqref{eq:lorentzian-connected-derivative-normalization} and
\eqref{eq:euclidean-connected-derivative-normalization}, together with the
analytic continuation of the connected distributions, are the matching rule
used by the later Euclidean path-integral and RG chapters.  For operators with
Lorentz indices, spinor indices, or time derivatives, the usual Wick-rotation
matrix for those indices must also be applied.
\end{definition}

In particular, the Euclidean free-energy convention \(W_E=-\log Z_E\) is
opposite to the connected-log convention \(\mathcal W_E=\log Z_E\).  The source
derivative of \(W_E\) has the sign displayed in
\eqref{eq:euclidean-connected-derivative-normalization}; this is the sign used
in the convex Legendre transform below.

\begin{theorem}[Euclidean convexity at finite regulator]
\label{thm:volume-ii-euclidean-convexity-finite-regulator}
Assume Assumption~\ref{ass:positive-euclidean-regulator-convexity}.  Define
\[
  K[J]=\log Z_E[J],
  \qquad
  W_E[J]=-K[J].
\]
Thus \(W_E\) is the Euclidean free energy.  It is the negative of the
connected-log functional \(K\).  Functional derivative signs are therefore fixed
by two independent choices: whether one differentiates \(K\) or \(-K\), and the
sign with which the source is placed in the exponent.  For the present source
term \(-\langle J,\phi\rangle\), one has
\(\delta W_E/\delta J=\langle\phi\rangle_J\).  For a source term
\(+\langle J,\mathcal O\rangle\), the same free-energy convention gives
\(\delta W/\delta J=-\langle\mathcal O\rangle_J\).
Then \(K\) is convex on \(\mathcal D_\Lambda\), and \(W_E\) is concave on
\(\mathcal D_\Lambda\).  The finite-regulator Euclidean effective action
\[
  \Gamma_E[\varphi]
  =
  \sup_{J\in\mathcal D_\Lambda}
  \left\{
    W_E[J]-\langle J,\varphi\rangle
  \right\}
\]
is a convex functional of \(\varphi\).  At a differentiability point where the
supremum is attained by a locally unique source \(J_\varphi\) satisfying
\[
  \varphi(x)
  =
  \left.
  \frac{\delta W_E[J]}{\delta J(x)}
  \right|_{J=J_\varphi},
\]
this definition reduces to the local Legendre transform
\[
  \Gamma_E[\varphi]
  =
  W_E[J_\varphi]-\langle J_\varphi,\varphi\rangle,
  \qquad
  \frac{\delta\Gamma_E[\varphi]}{\delta\varphi(x)}
  =
  -J_\varphi(x),
\]
and its second variation is nonnegative on the subspace where the
source-field map is nonsingular.
\end{theorem}

\begin{proof}
For \(J_1,J_2\in\mathcal D_\Lambda\) and \(0\le t\le1\), Holder's inequality
with respect to the positive measure \(\mu_\Lambda\) gives
\[
\begin{aligned}
  Z_E[tJ_1+(1-t)J_2]
  &=
  \int_{\mathcal C_\Lambda}
  \left(
    \ee^{-\langle J_1,\phi\rangle_\Lambda}
  \right)^t
  \left(
    \ee^{-\langle J_2,\phi\rangle_\Lambda}
  \right)^{1-t}
  \dd\mu_\Lambda(\phi)          \\
  &\le
  Z_E[J_1]^t Z_E[J_2]^{1-t}.
\end{aligned}
\]
Taking the logarithm yields
\[
  K[tJ_1+(1-t)J_2]
  \le
  tK[J_1]+(1-t)K[J_2],
\]
so \(K\) is convex.  Equivalently,
\[
  W_E[tJ_1+(1-t)J_2]
  \ge
  tW_E[J_1]+(1-t)W_E[J_2].
\]
For each fixed source \(J\), the map
\(\varphi\mapsto W_E[J]-\langle J,\varphi\rangle\) is affine.  The supremum of
affine functionals is convex, proving convexity of \(\Gamma_E\).  If the
supremum is attained at a locally unique differentiability point, its
stationarity condition is
\(\varphi=\delta W_E/\delta J|_{J=J_\varphi}\).  Differentiating
\(W_E[J_\varphi]-\langle J_\varphi,\varphi\rangle\) then cancels the
\(\delta J_\varphi\) terms and gives
\(\delta\Gamma_E/\delta\varphi=-J_\varphi\).  Convexity gives, in any
nonsingular local chart,
\[
  \int \dd^Dx\,\dd^Dy\,
  f(x)\Gamma_E^{(2)}(x,y)f(y)\ge0
\]
for every real test function \(f\) in the corresponding tangent subspace.
\end{proof}

\section{Legendre--Fenchel Form at Finite Regulator}

This section remains in the positive-measure setting of
Theorem~\ref{thm:volume-ii-euclidean-convexity-finite-regulator}.  It is not a
replacement for the perturbative 1PI definition used in dimensional
regularization or in any other regulator whose ``path integral'' is only a
formal algebraic prescription.  The two constructions agree only in a domain
where both are defined and the convex conjugate is differentiable with a
locally invertible source-field map.

The local inverse \(J_{\varphi}\) used above is a convenient notation.
Theorem~\ref{thm:volume-ii-euclidean-convexity-finite-regulator} identifies the
exact finite-regulator definition as a convex-conjugate formula.  At finite
regulator, the source space is a finite-dimensional vector space or a specified
space of test functions, depending on the regulator.  Its dual is
the space in which the regulated mean field \(\varphi\) lives, and the pairing is
\[
  \langle J,\varphi\rangle
  =
  \int \dd^Dx\,J(x)\varphi(x)
\]
with the corresponding finite sum on a lattice.

Let
\[
  K[J]=\log Z_E[J],
  \qquad
  W_E[J]=-K[J].
\]
Under the positivity assumption of
Theorem~\ref{thm:volume-ii-euclidean-convexity-finite-regulator}, \(K\) is
convex.  The exact Euclidean effective action is the convex conjugate with the
sign convention fixed by the source term \(-\langle J,\phi\rangle\):
\[
  \Gamma_E[\varphi]
  =
  \sup_J
  \left\{
    -K[J]-\langle J,\varphi\rangle
  \right\}
  =
  \sup_J
  \left\{
    W_E[J]-\langle J,\varphi\rangle
  \right\}.
\]
This formula covers cases in which the map \(J\mapsto\varphi_J\) has multiple
preimages or the supremum is attained by more than one source.
It is also meaningful at boundary points of the domain where the supremum is
represented by a limiting sequence of sources.
If \(K\) is differentiable and the supremum is attained at \(J_\varphi\), the
stationarity condition is
\[
  0
  =
  -\frac{\delta K[J_\varphi]}{\delta J(x)}
  -
  \varphi(x),
  \qquad
  \varphi(x)
  =
  \frac{\delta W_E[J_\varphi]}{\delta J(x)} .
\]
Thus the earlier Legendre transform is the differentiable local form of the
Legendre--Fenchel definition.

The subgradient statement is the coordinate-free replacement for the inverse
source map.  A source \(J\) maximizes the displayed variational expression at
\(\varphi\) precisely when
\[
  -\varphi
  \in
  \partial K[J],
\]
where \(\partial K[J]\) is the convex subdifferential of \(K\) at \(J\).  If
\(K\) is twice differentiable and its Hessian has a kernel at a source \(J\),
the source-field map is locally degenerate there.  The convex-conjugate
formula remains the definition.  It records the supporting hyperplanes of
\(K\): an admissible mean field \(\varphi\) is one for which \(-\varphi\) is a
subgradient of \(K\) at some source, or a limit of such subgradients at a
boundary point of the domain.

A finite-dimensional model displays the same logic without functional
analysis.  Let \(\phi\in\R\) and
\[
  Z_E(J)=\int_{\R}\dd\phi\,\ee^{-S(\phi)-J\phi}
  =
  \ee^{-W_E(J)}
\]
with \(S(\phi)\) a confining function.  The source-dependent mean is
\[
  \varphi(J)=\frac{\dd W_E}{\dd J}.
\]
In the small-fluctuation limit, \(W_E(J)\) is obtained by minimizing
\(S(\phi)+J\phi\) with respect to \(\phi\).  The corresponding
\(\Gamma_E(\varphi)\) is the convex Legendre transform.  If \(S(\phi)\) has
more than one local minimum, different ranges of \(J\) select different
stationary branches, and the exact Legendre transform produces a convex
function of \(\varphi\).  A perturbative expansion around a chosen branch is
valid where the quadratic fluctuation operator at that branch is positive.

The last clause is an analytic condition, not a matter of terminology.  Choose
a point \(\phi_0\) on a stationary branch and choose the source
\[
  J_0=-S'(\phi_0).
\]
The Euclidean background representation of the effective action is
\[
  \ee^{-\Gamma_E(\phi_0)}
  =
  \left.
  \int \dd\xi\,
  \exp\{-S(\phi_0+\xi)-J_0\xi\}
  \right|_{\langle\xi\rangle_{\phi_0,J_0}=0}.
\]
The exponent is obtained by subtracting the tangent line to \(S\) at
\(\phi_0\).  Its expansion begins
\[
  S(\phi_0+\xi)+J_0\xi
  =
  S(\phi_0)
  +
  \frac12 S''(\phi_0)\xi^2
  +
  O(\xi^3).
\]
If \(S''(\phi_0)<0\), the Gaussian integral around this point is not a
positive quadratic approximation to the finite-dimensional Euclidean
integral.  The exact \(\Gamma_E\) is then determined by the global supporting
line data encoded in the Legendre--Fenchel transform, while the coefficientwise
stationary-phase expansion around that branch computes only a local branch
functional.  This is the
finite-dimensional form of the distinction between the exact convex effective
action and a perturbative effective potential computed around a specified
saddle.

Figure~\ref{fig:volume-ii-branch-perturbation-condition} records the local
branch condition for the finite-dimensional model: perturbation theory around a
chosen saddle uses a positive quadratic part, while the exact
Legendre--Fenchel transform uses global supporting lines.

\begin{figure}[!ht]
  \centering
  \resizebox{0.94\textwidth}{!}{%
  \begin{tikzpicture}[>=Stealth,x=1cm,y=1cm,every node/.style={font=\small}]
    \tikzset{axis/.style={->,line width=0.75pt}, curve/.style={line width=1.05pt}}
    \begin{scope}
      \node[font=\normalsize] at (0,2.2) {branch and tangent};
      \draw[axis] (-2.2,-1.25)--(2.35,-1.25) node[right] {\(\phi\)};
      \draw[axis] (-1.85,-1.45)--(-1.85,1.55) node[above] {\(S(\phi)\)};
      \draw[blue!70!black,curve]
        plot[smooth,tension=0.82] coordinates {
          (-2.0,1.35) (-1.45,-0.55) (-0.5,0.42) (0.28,0.05) (0.95,-0.62) (1.75,1.25)
        };
      \coordinate (p) at (0.25,0.08);
      \fill (p) circle (2.2pt);
      \node[above right] at (p) {\(\phi_0\)};
      \draw[orange!85!black,dashed,line width=0.9pt] (-1.35,1.25)--(1.45,-0.65);
      \node[orange!85!black] at (0.95,-1.02) {\(J_0=-S'(\phi_0)\)};
      \node[orange!85!black] at (-1.05,-0.74) {tangent line};
    \end{scope}

    \begin{scope}[shift={(6.1,0)}]
      \node[font=\normalsize] at (0,2.2) {shifted exponent};
      \draw[axis] (-2.25,-1.25)--(2.35,-1.25) node[right] {\(\xi\)};
      \draw[axis] (-1.85,-1.45)--(-1.85,1.55)
        node[above] {\(S(\phi_0+\xi)+J_0\xi\)};
      \draw[purple!80!black,curve]
        plot[smooth,tension=0.82] coordinates {
          (-1.9,1.05) (-1.15,-0.82) (-0.35,-0.18) (0,0.0) (0.55,-0.38) (1.45,1.15)
        };
      \fill[red!75!black] (0,0) circle (2.2pt);
      \node[red!75!black,align=center] at (0.75,-0.88)
        {\(S''(\phi_0)<0\)\\no Gaussian saddle};
      \draw[gray!65,dashed] (-0.65,0.55)--(0.65,-0.55);
    \end{scope}
  \end{tikzpicture}
  }
  \caption{In the finite-dimensional model, the source subtracts the tangent
  line at the chosen mean field.  Perturbation theory around that point
  requires the shifted exponent to have a positive quadratic part.  The exact
  Legendre--Fenchel effective action instead uses the global supporting-line
  structure and is convex.}
  \label{fig:volume-ii-branch-perturbation-condition}
\end{figure}

\clearpage

Figure~\ref{fig:volume-ii-euclidean-convexity} displays the convex Euclidean
effective action obtained from the exact Legendre--Fenchel construction, in
contrast to the nonconvex local branch just described.

\begin{figure}[!ht]
  \centering
  \resizebox{0.95\textwidth}{!}{%
  \begin{tikzpicture}[>=Stealth,x=1cm,y=1cm,every node/.style={font=\small}]
    \tikzset{axis/.style={->,line width=0.75pt}, curve/.style={line width=1.05pt}}
    \begin{scope}
      \node[font=\normalsize] at (0,2.25) {Euclidean response};
      \draw[axis] (-2.2,-1.35)--(2.35,-1.35) node[right] {\(J\)};
      \draw[axis] (-1.9,-1.6)--(-1.9,1.55) node[above] {\(W_E[J]\)};
      \draw[blue!70!black,curve]
        plot[smooth,tension=0.7] coordinates {
          (-2.05,-1.05) (-1.35,-0.05) (-0.45,0.72) (0.25,0.92) (1.05,0.58) (2.05,-0.95)
        };
      \draw[green!55!black,curve] (-1.35,-0.05)--(0.92,0.63);
      \draw[orange!85!black,curve] (-0.15,0.9)--(1.8,-0.65);
      \node[blue!70!black] at (0.6,1.2) {concave \(W_E\)};
      \node[green!50!black] at (-0.6,-0.65) {slope \(\varphi_1\)};
      \node[orange!85!black] at (1.15,0.1) {slope \(\varphi_2\)};
    \end{scope}

    \draw[-{Latex[length=2.0mm]},line width=0.9pt] (2.75,0) -- (4.2,0);
    \node[above] at (3.48,0) {Legendre};

    \begin{scope}[shift={(6.2,0)}]
      \node[font=\normalsize] at (0,2.25) {effective action};
      \draw[axis] (-2.1,-1.35)--(2.35,-1.35) node[right] {\(\varphi\)};
      \draw[axis] (-1.75,-1.6)--(-1.75,1.55) node[above] {\(\Gamma_E[\varphi]\)};
      \draw[gray!60,curve]
        plot[smooth,tension=0.85] coordinates {
          (-1.65,1.3) (-1.1,-0.35) (-0.45,-0.6) (0.05,0.05) (0.7,-0.56) (1.45,-0.25) (2.0,1.2)
        };
      \draw[purple!80!black,curve]
        plot[smooth,tension=0.8] coordinates {
          (-1.65,1.15) (-1.05,-0.3) (-0.45,-0.62) (0.15,-0.62) (0.75,-0.6) (1.35,-0.15) (1.95,1.08)
        };
      \node[gray!60!black] at (0.35,0.82) {stationary branches};
      \node[purple!80!black] at (0.55,-1.0) {convex \(\Gamma_E\)};
    \end{scope}
  \end{tikzpicture}
  }
  \caption{For a positive Euclidean measure and the convention
  \(Z_E[J]=\ee^{-W_E[J]}\), the response functional \(W_E\) is concave and the
  effective action \(\Gamma_E\) is convex.  In a finite-dimensional
  double-well model, the exact Legendre transform connects the stable branches
  by the convex segment.}
  \label{fig:volume-ii-euclidean-convexity}
\end{figure}

The Euclidean convexity theorem is a statement about the exact
Legendre--Fenchel transform of a positive finite-measure source functional.  A
Wilsonian action at a finite cutoff scale, the regulated action that appears
in the Euclidean measure, the local formal perturbative 1PI effective action,
and a perturbative effective potential computed around one stationary branch
are distinct functionals with distinct definitions.  The next chapter uses the
local formal 1PI functional as the arena in which perturbative local
counterterms are identified and organized.
